% !TEX TS-program = xelatex
% !BIB TS-program = bibtex
\documentclass[12pt,letterpaper]{article}
\usepackage{style/dsc180reportstyle} % import dsc180reportstyle.sty

%%%%%%%%%%%%%%%%%%%%%%%%%%%%%%%%%%%%%%%%%%%%%%%%%%%%%%%%
%%%% Title and Authors
%%%%%%%%%%%%%%%%%%%%%%%%%%%%%%%%%%%%%%%%%%%%%%%%%%%%%%%%

\title{
Toward Accelerated LLM Inference:\\
Porting and Evaluating Diffusion-Based\\
Speculative Decoding on TPU
}

\author{Aaron Feng \\
  {\tt zhf004@ucsd.edu} \\\And
  Zhongyan Luo \\
  {\tt zhl203@ucsd.edu} \\\And
  Son Nguyen \\
  {\tt hpn005@ucsd.edu} \\\And
  Andy Huang \\
  {\tt anh031@ucsd.edu} \\\AND
  Yiming Zhao (TA) \\
  {\tt yiz279@ucsd.edu} \\\And
  Hao Zhang (Mentor) \\
  {\tt haz094@ucsd.edu} \\}

\begin{document}
\maketitle

%%%%%%%%%%%%%%%%%%%%%%%%%%%%%%%%%%%%%%%%%%%%%%%%%%%%%%%%
%%%% Abstract and Links
%%%%%%%%%%%%%%%%%%%%%%%%%%%%%%%%%%%%%%%%%%%%%%%%%%%%%%%%

\begin{abstract}
Large language model (LLM) inference is bottlenecked by sequential autoregressive decoding, which requires one forward pass per generated token. Speculative decoding mitigates this limitation by introducing a lightweight draft model whose proposed tokens are verified in parallel by the target model. DFlash is a recent method that replaces autoregressive drafting with a block diffusion model, enabling parallel generation of multiple draft tokens in a single forward pass and demonstrating substantial speedups on GPU.

In this work, we adapt DFlash to TPU systems by reimplementing the algorithm in JAX within the \texttt{tpu-inference} framework. We evaluate on the same math reasoning benchmarks as the original DFlash paper, in both a standalone setting and within the vLLM serving pipeline. Our standalone results match the draft acceptance quality reported on GPU, with remaining speed differences attributable primarily to hardware characteristics. Within the vLLM pipeline, DFlash achieves substantial acceleration over the autoregressive baseline and retains its proportional advantage over autoregressive drafting methods such as Eagle3. We further verify that speculative decoding preserves output quality relative to standard decoding. These results demonstrate that diffusion-based speculative decoding transfers effectively to TPU and maintains its practical benefits over autoregressive drafting.
\end{abstract}

\begin{center}
\begin{tabular}{r l}
Code: & \url{https://github.com/aaronzhfeng/tpu-spec-decode} \\
      & \url{https://github.com/aaronzhfeng/tpu-inference}
\end{tabular}
\end{center}


\clearpage

\maketoc
\clearpage

%%%%%%%%%%%%%%%%%%%%%%%%%%%%%%%%%%%%%%%%%%%%%%%%%%%%%%%%
%%%% Main Contents
%%%%%%%%%%%%%%%%%%%%%%%%%%%%%%%%%%%%%%%%%%%%%%%%%%%%%%%%

\section{Introduction}

Autoregressive large language models (LLMs) generate text one token at a time: each token depends on all previously generated tokens, requiring a separate forward pass per output token. This sequential bottleneck dominates inference latency, particularly for long-form generation tasks such as chain-of-thought reasoning, where models may produce hundreds or thousands of tokens. As LLMs are deployed in latency-sensitive applications, reducing this per-token cost has become a central systems challenge.

\textbf{Speculative decoding} \citep{leviathan2023fastinferencetransformersspeculative} addresses this bottleneck by introducing a two-stage process: a lightweight \emph{draft model} quickly proposes a sequence of candidate tokens, and the full \emph{target model} verifies them in a single batched forward pass. Correctly predicted tokens are accepted for free; the first incorrect token triggers a fallback to the target model's own prediction. In the best case, many tokens are accepted per verification step, amortizing the cost of the target model's forward pass. The key metric is the \emph{acceptance length} $\tau$: the average number of tokens accepted per speculative step. Higher $\tau$ means more tokens generated per target-model call, yielding greater speedup.

Most existing speculative decoding methods, including Eagle3 \citep{li2025eagle3scalinginferenceacceleration}, Medusa \citep{cai2024medusasimplellminference}, and prompt lookup decoding, use autoregressive draft models that generate candidate tokens sequentially. While the draft model is cheaper than the target, the sequential drafting still limits throughput and introduces error accumulation, capping practical speedups at 2--3$\times$.

\textbf{DFlash} \citep{chen2026dflashblockdiffusionflash} takes a fundamentally different approach: it replaces the autoregressive drafter with a lightweight \emph{block diffusion model} that generates an entire block of 16 draft tokens in a single forward pass. The draft model is conditioned on hidden-state features extracted from the target model's intermediate layers, enabling high-quality parallel drafting without sequential dependencies between draft positions. On GPU, DFlash achieves over 5$\times$ speedup, substantially outperforming Eagle3's 2$\times$.

However, DFlash has only been implemented for GPU using PyTorch. Google Cloud TPUs offer a compelling alternative platform: their architecture, optimized for large, dense matrix operations with high memory bandwidth, is naturally suited to the parallel block-prediction pattern of diffusion-style drafting. The \texttt{tpu-inference} framework \citep{tpuinference} provides TPU support for vLLM with speculative decoding infrastructure (currently supporting Eagle3 and n-gram methods), but lacks a DFlash implementation.

\textbf{Our contribution.} We implement DFlash in JAX within the \texttt{tpu-inference} framework, extending its speculative decoding infrastructure to support diffusion-based drafting on TPU systems. The implementation includes integration into the TPU speculative decoding manager as well as a standalone benchmark independent of the vLLM serving pipeline. We evaluate the system in both standalone and serving settings and compare against existing autoregressive drafting methods. Our results show that diffusion-based speculative decoding transfers effectively to TPU hardware and retains its practical advantages over autoregressive drafting. Beyond replication, this work establishes a foundation for investigating the broader interaction between speculative decoding, diffusion models, and TPU hardware, a direction that is particularly promising given TPU’s architectural affinity for dense, parallel computation.

\section{Background}

\subsection{Speculative Decoding}

In standard autoregressive decoding, generating $n$ tokens requires $n$ sequential forward passes through the target model. Speculative decoding \citep{leviathan2023fastinferencetransformersspeculative} reduces this by introducing a draft-then-verify paradigm. Given the current context, a draft model proposes $k$ candidate tokens. The target model then processes all $k$ tokens in a single forward pass (similar to prefilling), producing logits at each position. The verification step compares draft tokens against the target's predictions: consecutive matches from the first position are accepted, and the first mismatch triggers rejection. If $m$ tokens are accepted ($m \leq k$), then the target model also provides the correct token at position $m+1$, yielding $m+1$ new tokens from a single target forward pass.

The expected speedup depends on the acceptance length $\tau$ (average tokens accepted per step) and the relative cost of drafting versus target verification. Higher $\tau$ means fewer target forward passes per generated token.

\subsection{DFlash: Block Diffusion for Speculative Decoding}

DFlash \citep{chen2026dflashblockdiffusionflash} replaces the autoregressive drafter with a block diffusion model. The key components are:

\begin{enumerate}
    \item \textbf{Target feature extraction.} During each target model forward pass, hidden states from a subset of intermediate layers (e.g., layers 1, 9, 17, 25, 33 for Qwen3-4B) are concatenated and projected to produce a context vector that summarizes the target model's ``reasoning state.''
    \item \textbf{Parallel block drafting.} The draft model receives (a) the context vector from the target, and (b) a noise block: the embedding of the last accepted token followed by mask tokens at 15 unknown positions. In a single forward pass with \emph{non-causal} attention over the block, the draft model predicts logits at all 16 positions simultaneously.
    \item \textbf{Shared vocabulary.} The draft model has no embedding layer or language model head of its own. It reuses the target model's \texttt{embed\_tokens} for input and the target's LM head for computing output logits, ensuring draft and target operate in the same vocabulary space.
    \item \textbf{Verification and acceptance.} The full block of 16 tokens is verified by the target model in one forward pass. Consecutive matches are accepted, and the bonus token from the target extends the sequence.
\end{enumerate}

On GPU (H200), DFlash achieves $\tau=7.07$ and 5.53$\times$ average speedup on Qwen3-4B across four math reasoning datasets, substantially outperforming Eagle3's $\tau=3.40$ and 2.02$\times$ \citep{chen2026dflashblockdiffusionflash}.

\subsection{TPU Inference and vLLM}

vLLM \citep{kwon2023efficientmemorymanagementlarge} is a high-throughput LLM serving engine that provides continuous batching, PagedAttention for efficient KV cache management, and scheduling for concurrent requests. The \texttt{tpu-inference} plugin \citep{tpuinference} adapts vLLM to run on Google Cloud TPUs, implementing model execution and attention in JAX. It supports speculative decoding via a modular \texttt{SpeculativeDecodingManager} that coordinates draft proposal, target verification, and token acceptance. Prior to our work, only Eagle3 and n-gram drafting were supported.


\section{Methods}

Porting DFlash from GPU/PyTorch to TPU/JAX required addressing several architectural mismatches between the two platforms. We describe the three main technical challenges and our solutions.

\subsection{Dual KV Cache Architecture}

The GPU DFlash implementation uses PyTorch's \texttt{DynamicCache} with simple concatenation-based KV management for both the target and draft models. On TPU, the target model uses \emph{paged attention} with Pallas kernels (\texttt{ragged\_paged\_attention}), which stores KV entries in fixed-size pages managed by block tables. The draft model's non-causal attention pattern is incompatible with paged attention.

We implemented a \textit{dual-cache architecture}: the target model retains its paged KV cache (required for \texttt{ragged\_paged\_attention}), while the draft model uses static on-device JAX arrays with \texttt{dynamic\_update\_slice} writes and a \texttt{cache\_len} pointer for crop semantics. The draft model uses TPU's \texttt{flash\_attention} kernel with \texttt{causal=False}, matching DFlash's non-causal block diffusion design.

\subsection{Context Buffer Management}

DFlash's draft model is conditioned on projected hidden states from the target model, accumulated across iterations in a \emph{context buffer}. The GPU implementation manages this in host-side numpy arrays with power-of-2 padding for efficient TPU transfers. We replicated this pattern: after each target model forward pass, newly projected features are appended to the context buffer, padded to the next power of two, and passed to the draft model alongside a \texttt{cache\_len} and \texttt{actual\_ctx\_count} to indicate which features are new. This required careful tracking of \texttt{prev\_ctx\_len} (how much context the draft model has already consumed) to avoid feeding duplicate or missing context.

\subsection{Sequence Length Mismatch in Speculative Verification}

During integration into the vLLM pipeline, initial acceptance rates were significantly lower than expected ($\tau=2.49$, 1.30$\times$ speedup). Investigation revealed that the speculative decoding manager forwarded \texttt{attn\_metadata.seq\_lens} to the proposer during verification steps. In this setting, \texttt{seq\_lens} reflects the total sequence length including draft tokens under verification, rather than the number of tokens actually accepted.

Because DFlash maintains persistent state across iterations (including context buffers, KV cache positions, and RoPE offsets), this discrepancy caused systematic misalignment between internal draft state and the true accepted token count. The resulting sequence length inflation introduced cumulative state errors across iterations, degrading acceptance length and overall speedup.

Correcting the proposer to use the true accepted token count (\texttt{num\_tokens\_no\_spec}) restored alignment between draft state and target verification. This change increased acceptance length from $\tau=2.49$ to 4.48 and improved end-to-end speedup from 1.30$\times$ to 2.31$\times$.

\subsection{Standalone Benchmark}

To provide a direct comparison with the original DFlash results on GPU (which report standalone numbers without serving overhead) and to isolate the implementation from serving-layer effects, we built a standalone JAX benchmark (\texttt{benchmarks/standalone\_dflash.py}). This script loads models directly via \texttt{get\_flax\_model()}, runs a prefill-draft-verify-accept loop, and tracks context buffers and KV state manually, using only existing \texttt{tpu-inference} library imports (\texttt{ModelConfig}, \texttt{LoadConfig}, model loader). No vLLM engine, scheduler, or serving pipeline is involved.



\section{Results}

We evaluate on four math reasoning datasets used in the DFlash paper: GSM8K, Math500, AIME24, and AIME25. All experiments use Qwen3-4B as the target model and z-lab/Qwen3-4B-DFlash-b16 as the draft model, with greedy decoding (temperature = 0), block size = 16, and 8 samples per dataset (7 measured, 1 warmup for JIT compilation).

\subsection{Standalone Benchmark: TPU vs.\ Original DFlash (GPU)}

Table \ref{tab:standalone} shows our standalone TPU results compared with the original DFlash numbers on GPU.

\begin{table}[htbp]
\caption{Standalone benchmark results: TPU (single v4 device, JAX) vs.\ original DFlash on GPU (H200, PyTorch). $\tau$ = average acceptance length, Speedup = wall-clock ratio over autoregressive baseline, TPS = tokens per second.}
\label{tab:standalone}
\resizebox{\linewidth}{!}{\centering
\begin{tabular}{l|cccc|cc}
\toprule
Dataset & TPU $\tau$ & TPU Speedup & TPU TPS & Baseline TPS & GPU $\tau$ & GPU Speedup \\
\midrule
GSM8K   & 5.17 & 3.36$\times$ & 312.2 & 93.0 & 6.53 & 5.15$\times$ \\
Math500 & \textbf{8.72} & \textbf{5.70}$\times$ & 531.5 & 93.2 & 7.84 & 6.09$\times$ \\
AIME24  & 6.32 & 3.93$\times$ & 367.7 & 93.6 & 7.27 & 5.68$\times$ \\
AIME25  & 6.48 & 4.18$\times$ & 388.9 & 93.1 & 6.64 & 5.21$\times$ \\
\midrule
\textbf{Average} & \textbf{6.67} & \textbf{4.29}$\times$ & \textbf{400.1} & 93.2 & 7.07 & 5.53$\times$ \\
\bottomrule
\end{tabular}
}
\end{table}

Our TPU implementation achieves $\tau=6.67$ on average, which is 94\% of the original DFlash $\tau=7.07$ on GPU. On Math500, our TPU $\tau$ of 8.72 exceeds the original 7.84. The average end-to-end speedup is 4.29$\times$ compared to the original 5.53$\times$ on GPU. The remaining speedup gap is primarily due to hardware differences (TPU v4 vs.\ H200 GPU) rather than algorithm quality, as the $\tau$ ratio is much closer than the speedup ratio.

\subsection{vLLM Pipeline Integration}

We also integrated DFlash into the full vLLM serving pipeline on TPU (4 TPU v4 devices). Table \ref{tab:cross} shows the cross-comparison.

\begin{table}[htbp]
\caption{Cross-comparison: standalone TPU, vLLM pipeline TPU, and original DFlash (GPU). The vLLM pipeline shows lower speedup due to scheduling and batch management overhead, not algorithm quality.}
\label{tab:cross}
\resizebox{\linewidth}{!}{\centering
\begin{tabular}{l|cc|cc|cc}
\toprule
 & \multicolumn{2}{c|}{Standalone (TPU)} & \multicolumn{2}{c|}{vLLM Pipeline (TPU)} & \multicolumn{2}{c}{Original DFlash (GPU)} \\
Dataset & $\tau$ & Speedup & $\tau$ & Speedup & $\tau$ & Speedup \\
\midrule
GSM8K   & 5.17 & 3.36$\times$ & --- & 2.16$\times$ & 6.53 & 5.15$\times$ \\
Math500 & 8.72 & 5.70$\times$ & --- & 2.62$\times$ & 7.84 & 6.09$\times$ \\
AIME24  & 6.32 & 3.93$\times$ & --- & 2.09$\times$ & 7.27 & 5.68$\times$ \\
AIME25  & 6.48 & 4.18$\times$ & --- & 2.36$\times$ & 6.64 & 5.21$\times$ \\
\midrule
\textbf{Average} & \textbf{6.67} & \textbf{4.29}$\times$ & \textbf{4.48} & \textbf{2.31}$\times$ & \textbf{7.07} & \textbf{5.53}$\times$ \\
\bottomrule
\end{tabular}
}
\end{table}

The vLLM pipeline achieves $\tau=4.48$ and 2.31$\times$ speedup (212 TPS vs.\ 92 TPS baseline). The lower numbers compared to the standalone benchmark are expected: vLLM's scheduler, batch manager, and rejection sampling loop add per-step overhead that compounds across speculative decoding iterations. This matches the GPU side, where the paper's standalone results are also significantly better than their serving numbers. Furthermore, vLLM's TPU v4 support is listed as experimental; we plan to rerun on v5 TPU (where vLLM is fully supported) once our requested instance becomes available.

\subsection{Comparison with Eagle3}

The only other speculative decoding method available on TPU through \texttt{tpu-inference} is Eagle3. Table \ref{tab:eagle3} compares DFlash and Eagle3 on the same vLLM pipeline.

\begin{table}[htbp]
\caption{DFlash vs.\ Eagle3 on the vLLM TPU pipeline. Different target models (Qwen3-4B vs.\ Llama-3.1-8B), so not a direct head-to-head, but shows the advantage of block diffusion (15 tokens/step) over autoregressive drafting (3 tokens/step).}
\label{tab:eagle3}
\resizebox{0.85\linewidth}{!}{\centering
\begin{tabular}{l|ccc|ccc}
\toprule
 & \multicolumn{3}{c|}{DFlash (Qwen3-4B)} & \multicolumn{3}{c}{Eagle3 (Llama-3.1-8B)} \\
Metric & Value & & & Value & & \\
\midrule
Draft tokens/step & 15 & & & 3 & & \\
$\tau$ (acceptance length) & 4.48 & & & 2.18 & & \\
Per-token acceptance & 23.2\% & & & 39.3\% & & \\
Tokens/sec & 212 & & & 78 & & \\
Baseline TPS & 92 & & & 51 & & \\
Speedup & 2.31$\times$ & & & 1.53$\times$ & & \\
\bottomrule
\end{tabular}
}
\end{table}

DFlash's $\tau$ is 2.06$\times$ that of Eagle3 (4.48 vs.\ 2.18) and its speedup is 1.51$\times$ that of Eagle3 (2.31$\times$ vs.\ 1.53$\times$). These ratios closely match the original DFlash proportional advantage on GPU: DFlash's $\tau$ is 1.91$\times$ Eagle3's (6.49 vs.\ 3.40). This suggests DFlash's fundamental advantage over autoregressive drafting transfers well across hardware platforms.

\subsection{Output Quality Verification}

Speculative decoding with greedy sampling is mathematically guaranteed to produce the same output as autoregressive decoding under exact arithmetic. In practice, bf16 numerical precision can cause minor divergences. We verified output quality by comparing DFlash generations against the autoregressive baseline token-by-token on GSM8K (8 samples, 1024 max tokens).

\begin{table}[htbp]
\caption{Output quality check on GSM8K. 4/8 samples produce bit-exact identical output; 7/8 arrive at the same final $\backslash$\texttt{boxed\{\}} math answer. Divergences occur at random positions where top-2 logits are close, consistent with bf16 precision effects.}
\label{tab:quality}
\resizebox{\linewidth}{!}{\centering
\begin{tabular}{l|cccccc}
\toprule
Sample & Token Match & First Mismatch & Mismatched Tokens & Output Length & Baseline Ans. & DFlash Ans. \\
\midrule
0 & Yes & --- & 0 & 288 & 109 & 109 \\
1 & No & pos 245 & 404 & 681 & 92 & 92 \\
2 & No & pos 8 & 196 & 247 & 13 & 13 \\
3 & Yes & --- & 0 & 239 & 5 & 5 \\
4 & No & pos 13 & 276 & 351 & 25 & 25 \\
5 & No & pos 68 & 453 & 532 & 452 & 470 \\
6 & Yes & --- & 0 & 273 & 43 & 43 \\
7 & Yes & --- & 0 & 258 & 34 & 34 \\
\midrule
\multicolumn{2}{l}{\textbf{Exact token match:}} & \multicolumn{2}{l}{4/8 (50\%)} & \multicolumn{2}{l}{\textbf{Same answer:}} & 7/8 (87.5\%) \\
\bottomrule
\end{tabular}
}
\end{table}

As shown in Table \ref{tab:quality}, 4 out of 8 samples produce bit-exact identical output. The remaining 4 samples diverge at a single token position (positions 8, 13, 68, and 245), after which subsequent tokens cascade differently due to autoregressive dependencies. Despite these divergences, 7 out of 8 samples arrive at the same final boxed answer. The one answer difference (452 vs.\ 470) stems from an early divergence at token 68 where the model chooses a different number, leading to a different arithmetic path. The divergence points show coherent alternative text (e.g., \texttt{**Indras**} vs.\ \texttt{Indras}, \texttt{Step 1:} vs.\ \texttt{Initial Setup:}), confirming these are decision-boundary effects, not corruption. DFlash does not degrade model intelligence.


\section{Discussion}

\subsection{Why Is the Standalone Faster Than vLLM?}

The standalone benchmark achieves 4.29$\times$ speedup while the vLLM pipeline achieves 2.31$\times$, despite using the same model and algorithm. This is because vLLM is a \emph{serving} framework designed for high-throughput concurrent request handling, not single-request latency. Its scheduler, batch manager, PagedAttention bookkeeping, and rejection sampling loop add per-step overhead. Speculative decoding amplifies this overhead because it performs more round trips (draft + verify) per generated token than standard decoding. The original DFlash results show the same pattern: standalone numbers are substantially better than their SGLang serving numbers.

Notably, our vLLM benchmark runs requests sequentially (\texttt{max\_num\_seqs=1}), which represents the worst case for vLLM: all scheduling overhead with no batching benefit. The original DFlash serving benchmarks use concurrent requests (concurrency 1--32) with a ThreadPoolExecutor. A fairer serving comparison on TPU would require running vLLM in online serving mode with concurrent clients, which we leave to future work.

\subsection{TPU vs.\ GPU: Where Does the Gap Come From?}

Our $\tau$ (6.67) is 94\% of the original DFlash $\tau$ on GPU (7.07), but our speedup (4.29$\times$) is 78\% of the original (5.53$\times$). The $\tau$ gap is small and may be attributable to minor differences in model weight precision or attention kernel numerics. The larger speedup gap likely comes from hardware characteristics: the H200 GPU has higher memory bandwidth and may have faster kernel dispatch for the specific computation patterns in speculative decoding (repeated small matrix operations with varying shapes).

\subsection{vLLM Independence}

A key design goal was ensuring the DFlash implementation has zero dependency on vLLM changes. The standalone benchmark imports only \texttt{ModelConfig} and \texttt{LoadConfig} from vLLM, which are existing \texttt{tpu-inference} dependencies used by every model. No new vLLM code is required. The vLLM serving integration requires only a $\sim$10-line config change in a separate repository, which can be upstreamed independently. This means a pull request to \texttt{tpu-inference} can proceed without any coordination with the vLLM project.


\section{Future Work}

With the DFlash TPU replication validated, we see two categories of next steps: near-term engineering work and longer-term research directions at the intersection of speculative decoding, diffusion models, and TPU hardware.

\subsection{Near-Term Engineering}

\begin{enumerate}
    \item \textbf{Pull request to \texttt{tpu-inference}.} Prepare and submit the DFlash implementation (model, proposer, standalone benchmark, seq\_len fix) as a contribution to the upstream repository.
    \item \textbf{TPU v5 benchmarks.} Our current results are on TPU v4, where vLLM support is listed as experimental. We have requested a v5 instance and plan to rerun all benchmarks, as vLLM lists v5 as fully supported.
    \item \textbf{Concurrent serving benchmarks.} Run the vLLM pipeline with multiple concurrent clients at various concurrency levels (1, 4, 8, 16, 32) to match the original DFlash serving evaluation methodology and measure throughput scaling.
    \item \textbf{Multi-device standalone.} Add \texttt{shard\_map} wrapping for the \texttt{flash\_attention} Pallas kernel to enable multi-device standalone runs.
\end{enumerate}

\subsection{Speculative Decoding and Diffusion Models on TPU}

Our DFlash port raises a broader research question: \textit{how do speculative decoding methods interact with diffusion-based generation on TPU?} To our knowledge, no prior work has benchmarked diffusion-based speculative decoding on TPU, making our DFlash port the first data point in this space.

DFlash uses a diffusion model as a \emph{drafter} for an autoregressive target, but diffusion language models (dLLMs) are increasingly used as the \emph{target} model itself. Recent work such as FailFast \citep{pan2026failfastwinbig} explores speculative decoding where both the drafter and verifier operate in the diffusion paradigm, with adaptive speculation lengths based on drafter confidence. This is complementary to our work: while DFlash accelerates autoregressive targets with diffusion drafting, FailFast accelerates diffusion targets with speculative verification.

One hypothesis worth investigating is whether TPU hardware is particularly well-suited to diffusion-based speculative decoding. Diffusion models rely on dense, parallel matrix operations over blocks of tokens rather than sequential token-by-token generation. TPU systolic arrays are designed for such data-parallel workloads. Recent work on compute-in-memory TPU architectures has shown that compute-bound workloads (such as image Diffusion Transformers) benefit more from TPU scaling than memory-bound workloads like autoregressive decoding \citep{wang2025acceleratinglargelanguagemodel}. Whether this advantage extends to discrete diffusion language models remains an open question that has not been empirically tested.

We plan to investigate these directions by: (1) benchmarking diffusion language models (e.g., Fast-dLLM) as both drafters and targets on TPU, (2) evaluating whether the FailFast adaptive speculation strategy transfers to TPU hardware, and (3) comparing TPU vs.\ GPU speedups for diffusion-based vs.\ autoregressive-based speculative decoding to test whether the hypothesized TPU advantage holds in practice.

\section{Conclusion}

We implemented DFlash speculative decoding in JAX within the \texttt{tpu-inference} framework, extending TPU support to diffusion-based drafting. In standalone benchmarks on a single TPU v4 device, our implementation achieves an average 4.29$\times$ end-to-end speedup over autoregressive decoding while reaching 94\% of the draft acceptance length reported on GPU. Within the vLLM serving pipeline, the system achieves 2.31$\times$ speedup over the autoregressive baseline without degrading output quality.

These results demonstrate that diffusion-based speculative decoding transfers effectively from GPU to TPU and retains its proportional advantages over autoregressive drafting in practice. Beyond reproducing prior results on a new hardware platform, this work establishes a foundation for investigating the interaction between speculative decoding, diffusion models, and TPU architectures, particularly in settings where dense, parallel computation aligns with TPU design.

\clearpage

%%%%%%%%%%%%%%%%%%%%%%%%%%%%%%%%%%%%%%%%%%%%%%%%%%%%%%%%
%%%% Reference / Bibliography
%%%%%%%%%%%%%%%%%%%%%%%%%%%%%%%%%%%%%%%%%%%%%%%%%%%%%%%%

\makereference

\bibliography{reference}
\bibliographystyle{style/dsc180bibstyle}

%%%%%%%%%%%%%%%%%%%%%%%%%%%%%%%%%%%%%%%%%%%%%%%%%%%%%%%%
%%%% Appendix
%%%%%%%%%%%%%%%%%%%%%%%%%%%%%%%%%%%%%%%%%%%%%%%%%%%%%%%%

\clearpage
\makeappendix

\subsection{Files Modified in \texttt{tpu-inference}}

\begin{table}[htbp]
\caption{Files added or modified for the DFlash TPU port.}
\centering
\begin{tabular}{l|l}
\toprule
File & Change \\
\midrule
\texttt{models/jax/dflash.py} & DFlash model (new) \\
\texttt{spec\_decode/jax/dflash.py} & DFlash proposer (new) \\
\texttt{runner/speculative\_decoding\_manager.py} & seq\_len fix + DFlash integration \\
\texttt{benchmarks/standalone\_dflash.py} & Standalone benchmark (new) \\
\texttt{tests/standalone\_benchmark.sh} & Docker wrapper (new) \\
\bottomrule
\end{tabular}
\end{table}

\subsection{Reproduction Commands in \texttt{tpu-spec-decode}}

\begin{verbatim}
# Standalone benchmark (single dataset)
DATASET=math500 bash tests/standalone_benchmark.sh \
  --max-samples 8 --max-new-tokens 256

# All four datasets
for ds in gsm8k math500 aime24 aime25; do
  DATASET=$ds bash tests/standalone_benchmark.sh \
    --max-samples 8 --max-new-tokens 256
done

# vLLM pipeline benchmark
bash tests/benchmark.sh math
\end{verbatim}

\subsection{Latest Project Proposal}

\subsubsection*{Project Motivation}

Large language model (LLM) inference remains constrained by the sequential nature of autoregressive decoding. Speculative decoding mitigates this bottleneck by using a lightweight draft model whose proposed tokens are verified in parallel by a target model. Recent work (DFlash) replaces autoregressive drafting with a block diffusion model, enabling parallel generation of multiple draft tokens and achieving significant speedups on GPU.

However, diffusion-based speculative decoding has not been implemented or evaluated on TPU systems. Given TPU’s architectural affinity for dense, parallel matrix operations, diffusion-based drafting may align particularly well with TPU hardware. This project aims to port DFlash to TPU within the \texttt{tpu-inference} framework, evaluate its performance relative to GPU results, and compare it against existing autoregressive drafting methods such as Eagle3.

\subsubsection*{Project Objectives}

\begin{enumerate}
    \item Implement DFlash in JAX within the \texttt{tpu-inference} framework.
    \item Integrate DFlash into the TPU speculative decoding manager and vLLM serving pipeline.
    \item Build a standalone benchmark to isolate algorithmic performance from serving-layer overhead.
    \item Evaluate draft acceptance length, end-to-end speedup, and output quality on math reasoning benchmarks.
    \item Compare diffusion-based drafting against autoregressive drafting (Eagle3) on TPU.
\end{enumerate}

\subsubsection*{Six-Week Plan}

\textbf{Week 1: Background Study and System Design}

\begin{itemize}
    \item Study the DFlash paper and GPU implementation.
    \item Analyze TPU architecture and the \texttt{tpu-inference} codebase.
    \item Identify architectural differences between GPU/PyTorch and TPU/JAX execution.
    \item Define integration points within speculative decoding infrastructure.
\end{itemize}

\textbf{Week 2: Core DFlash Port to JAX}

\begin{itemize}
    \item Implement the DFlash draft model in JAX.
    \item Reproduce block diffusion drafting logic with non-causal attention.
    \item Implement shared embedding and LM head reuse.
    \item Validate functional correctness on small test cases.
\end{itemize}

\textbf{Week 3: KV Cache Integration and Initial Benchmarking}

\begin{itemize}
    \item Design dual KV cache architecture for TPU.
    \item Implement context buffer accumulation and feature projection.
    \item Integrate proposer into the speculative decoding manager.
    \item Build standalone JAX benchmark and run initial experiments.
    \item Measure draft acceptance length ($\tau$), tokens per second, and preliminary speedup.
\end{itemize}

\textbf{Week 4: vLLM Integration and Comparative Evaluation}

\begin{itemize}
    \item Integrate DFlash into the full vLLM TPU serving pipeline.
    \item Benchmark serving-mode performance.
    \item Compare against autoregressive drafting (Eagle3) on TPU.
    \item Verify output quality relative to autoregressive decoding.
\end{itemize}

\textbf{Week 5: Extended Benchmarking and Analysis}

\begin{itemize}
    \item Expand experiments across GSM8K, Math500, AIME24, and AIME25.
    \item Analyze TPU vs.\ GPU performance gaps.
    \item Investigate serving vs.\ standalone performance differences.
    \item Consolidate benchmark results for reporting.
\end{itemize}

\textbf{Week 6: Documentation, Website, and Final Deliverables}

\begin{itemize}
    \item Finalize reproducibility scripts and repository structure.
    \item Prepare structured project website summarizing motivation, methods, benchmarks, and findings.
    \item Draft final report and presentation materials.
    \item Prepare code for potential upstream contribution to \texttt{tpu-inference}.
\end{itemize}


\subsubsection*{Team Roles}

\begin{itemize}
    \item \textbf{Aaron Feng:} Led integration of DFlash into the \texttt{tpu-inference} framework, including proposer wiring, speculative decoding manager integration, TPU environment setup (v4/v5p), and serving-pipeline validation. Contributed to performance debugging and end-to-end system testing.

    \item \textbf{Zhongyan Luo:} Implemented the JAX standalone version of DFlash, including draft model execution, block diffusion logic, and prefill–draft–verify–accept loop. Led standalone benchmarking, draft acceptance measurement, and comparative evaluation against baseline decoding.

    \item \textbf{Son Nguyen:} Developed the conceptual framing of the project, designed evaluation methodology, analyzed TPU vs.\ GPU performance characteristics, and led report writing, documentation, and project website development. Contributed targeted integration fixes and validation of output correctness.

    \item \textbf{Andy Huang:} Assisted with speculative decoding integration support, experimental configuration management, benchmark orchestration, and result aggregation to ensure reproducibility.
\end{itemize}

\subsubsection*{Expected Outcomes}

By the end of six weeks, we expect to:

\begin{itemize}
    \item Deliver a working TPU implementation of DFlash integrated into \texttt{tpu-inference}.
    \item Provide standalone and serving-mode benchmarks on TPU.
    \item Demonstrate whether diffusion-based speculative decoding retains its proportional advantages over autoregressive drafting on TPU.
    \item Release reproducible code and documentation for future research.
\end{itemize}

\clearpage

\section*{Author Contributions}
\addcontentsline{toc}{section}{Author Contributions}

\textbf{Aaron Feng}

Aaron led the TPU integration of DFlash within the \texttt{tpu-inference} framework. He implemented proposer wiring, extended the speculative decoding manager to support diffusion-based drafting, and resolved the sequence-length alignment bug that affected acceptance rates. He configured TPU environments (v4 and v5p), validated serving-pipeline execution, and ensured compatibility with vLLM’s speculative decoding infrastructure. He played a central role in debugging integration issues, stabilizing the end-to-end TPU decoding stack, and preparing the implementation for upstream contribution.

\medskip

\textbf{Zhongyan Luo}

Zhongyan implemented the standalone JAX version of DFlash, including the block diffusion drafting logic, non-causal attention configuration, and the full prefill–draft–verify–accept decoding loop. He designed and validated the dual KV cache handling for standalone execution, ensuring correct context buffer updates and draft–target state alignment. He led standalone benchmarking, measured draft acceptance length and throughput, and conducted comparative evaluation against autoregressive baselines. He also contributed to debugging performance discrepancies and validating numerical consistency across runs.


\medskip

\textbf{Son Nguyen}

Son developed the conceptual framing of the project and structured the evaluation methodology. He analyzed TPU versus GPU performance differences, interpreted benchmarking results, and led report writing and documentation. He contributed integration validation, correctness verification, and preparation of the project website and reproducibility materials. He additionally explored comparative integration of LLaMA-based DFlash variants within the TPU framework.

\medskip

\textbf{Andy Huang}

Andy contributed to experimental configuration management and benchmark orchestration across standalone and serving settings. He assisted in coordinating dataset runs, organizing evaluation scripts, and aggregating results for comparison. He supported integration testing by verifying execution stability across TPU configurations and checking metric consistency across experiments. He also contributed to documentation refinement and preparation of final project deliverables.

\end{document}
